\documentclass[12pt]{article}
\usepackage[utf8]{inputenc}
\usepackage[T1]{fontenc}
\usepackage[brazil]{babel}
\usepackage[a4paper,margin=2.5cm]{geometry}
\usepackage{enumitem}
\setlist[itemize]{noitemsep,topsep=4pt}
\title{Sprint 3 - S3\_13\\Lógica Proposicional}
\author{Matheus Moreli Durão Silva\\Matrícula 12621CBS217\\Cibersegurança - Bacharelado - Noturno}
\date{2026}
\begin{document}
\maketitle

\section*{Entregável}
Lista resolvida de tabelas-verdade e argumentos. Este documento registra o objetivo do card, a solução desenvolvida e a forma de conferência. Os arquivos de código e as evidências citadas permanecem na mesma pasta deste documento.

\section*{O que foi estudado e aplicado}
\begin{itemize}
\item Conectivos
\item Tabelas-verdade
\item Leis de De Morgan
\item Argumentos válidos e inválidos
\end{itemize}

\section*{Como conferir}
O avaliador pode abrir os arquivos desta pasta, executar os exemplos quando houver código e comparar a saída com a explicação apresentada. Nos cards ligados ao Inventário de Ativos, a solução foi relacionada ao Trabalho 2 sem modificar requisitos já atendidos do CRUD.

\section*{Tabela e equivalências}
Com $P$ = ativo existente e $Q$ = CVSS válido, o cadastro exige $P\land Q$. Para $(P,Q)=(F,F),(F,V),(V,F),(V,V)$, $P\land Q$ resulta $F,F,F,V$; a implicação $P\to Q$ resulta $V,V,F,V$. As colunas de $\neg(P\land Q)$ e $\neg P\lor\neg Q$ são iguais: $V,V,V,F$, demonstrando De Morgan. $P\lor\neg P$ é tautologia e $P\land\neg P$ é contradição.
\section*{Argumentos}
Modus ponens é válido: de $P\to Q$ e $P$ segue $Q$, sem linha com premissas verdadeiras e conclusão falsa. A afirmação do consequente é inválida: de $P\to Q$ e $Q$ não segue $P$; $(P,Q)=(F,V)$ é contraexemplo. A lista \texttt{resolucao\_logica.md} traz tabelas completas, modus tollens e aplicação à validação do inventário. O script confirma as equivalências.
\section*{Conclusão}
Formalizar premissas evita confundir uma condição suficiente com necessária em regras de software.
\end{document}
